

\chapter{Skills}

\section{Writing}
Players can create notes at any time for their own rememberance, 
but these are not in-game items that can be passed on to new characters.

\vspace{0.4em}

\textit{Writing} exists for the easy transferral of information between characters,
and mostly consists of collecting useful knowledge into an in-game item,
but also can record unique encounters, spells, potions, etc.

\vspace{0.4em}

When a player levels-up in a town with a Library, they may choose 1:
\begin{itemize}
  \item if the player can \textit{Inscribe}, they may create 2 folios
  \item if the player can \textit{Map}, they may create 1 map.
\end{itemize}


\subsection{Tomes}
\textit{Tomes} are comprised of a variable number of \textit{Folios} (pages) that can contain information.
Tomes are created with a number of blank folios, which may be inscribed, 
and may have folios added to them, up to the maximum number of folios.

\vspace{0.4em}

The Storage Class of a Tome depends on the maximum number of folios it is able to contain.

\begin{Table}{c c}
  \TableHeader{\# Folios} & \TableHeader{Storage Class} \\
  < 10 & Small \\
  10--20 & Medium \\
  > 10 & Large \\
\end{Table}

The number of Folios required to record information depends on the type of information:
\begin{Table}{l c X}
  \TableHeader{Type} & \TableHeader{\# Folios} & \TableHeader{Storage Class} \\
  Spell & 2 & Teach the procedure for casting a spell. \\
  Potion & 1 & Teach the ingredients \& procedure for crafting a potion. \\
  Fact & 1 & Record an important plot point or fact of lore, and all essential context to understand that fact. \\
  Portrait & 1 & Visually depict a character, item, flora or fauna in an illustration. \\
  Landscape & 1 & Visually depict a landscape from a vantage point you have visited. \\
  Map & 2 & Record all spatial facts a player would require to navigate the depicted region. \\
\end{Table}



\section{Studying}
\textit{Studying Lore} allows players to safely gain knowledge about the world.
If a character can study lore, 
they may choose one of the following effects whenever they level-up in a town with a Library.

\begin{Table}{l X}
  \TableHeader{Name} & \TableHeader{Effect} \\
  Reflection & Roll \Roll{1d20} and add your \textbf{Intelligence} modifier. The DM tells the player about any narrative clues you missed since the last time you leveled-up. \\
  Preparation & Gain advantage on all Story (\Roll{d20}) rolls for \textbf{Wisdom \& Intelligence} until the next time you level-up. \\
  Exploration & Roll \Roll{1d20} and add your \textbf{Wisdom} modifier. The DM tells the player a narrative clue about an upcoming mission. \\
  Edification & Choose a topic and roll \Roll{1d20}, adding your \textbf{Intelligence} modifier. The DM tells the player information about the chosen topic. \\
\end{Table}



\section{Arcana} 
\textit{Arcana} allows characters to create new spells by synthesizing two spells into a single effect.

\vspace{0.4em}

If a character can use Arcana, they may craft one spell whenever they 
level-up in a town with an Arcaneum.

\vspace{0.4em}

To craft a new spell, choose two spells that you already know and wish to combine. 
Roll a \Roll{d10} to determine the outcome:

\begin{Table}{c X}
  \TableHeader{Roll} & \TableHeader{Outcome} \\
  1--2 & You immediately suffer the effects of both spells as if you were the target of those spells. \\
  3--4 & Nothing happens. \\
  5--8 & Gain advantage the next time you craft the same two spells. \\
  9--10 & Combine the effects and costs of the spells into a new spell, and give it a name. Its range becomes the lesser range of the two spells. You learn that spell. \\
\end{Table}


\newpage


\section{Crafting}
\textit{Crafting} allows characters to create or improve weapons and armor.

\vspace{0.4em}

If a character can craft weapons and armor, 
they may craft or improve one weapon or armor piece whenever they level-up in a town with a Workshop.

\subsection{Crafting Cost}
When you create or improve a weapon, you must expend a number of metal ingots based on the size:
\begin{Table}{c c}
  \TableHeader{\# Hands} & \TableHeader{\# Ingots} \\
  1 & 1 \\
  1/2 & 2 \\
  2 & 3 \\
\end{Table}

When you create or improve armor, 
you must expend a number of metal ingots and/or textiles based on the Armor Class:
\begin{Table}{l c c}
  \TableHeader{Armor Class} & \TableHeader{\# Ingots} & \TableHeader{\# Textiles} \\
  Unarmored & 0 & 2 \\
  Light & 1 & 1 \\
  Heavy & 2 & 0 \\
\end{Table}


\subsection{Improvement}
\textit{Improvements} are effects that are added to weapons or armor through crafting 
or enchantment, and are often campaign-specific.

\vspace{0.4em}

Standard weapons and armor may have up to \textbf{2} improvements.
Each item of equipment can only have one improvement with a given name.

\vspace{0.4em}

These improvements can be added to a weapon (other than a Mana Focus) in any campaign:

\begin{Table}{l X}
  \TableHeader{Improvement} & \TableHeader{Effect} \\
  Intricate & Add +2 to the maximum number of improvements for this weapon \\
  Lightweight & All W.I.Z.A.R.D.S. requirements for this weapon can be satisfied with Dexterity \\
  Honed & Add +1 to all Damage rolls for attacks using this weapon \\
\end{Table}

These improvements can be added to armor in any campaign:

\begin{Table}{l X}
  \TableHeader{Improvement} & \TableHeader{Effect} \\
  Intricate & Add +2 to the maximum number of improvements for this armor \\
  Fitted & Reduce the W.I.Z.A.R.D.S. attribute penalty by 1 (to a minimum of 0) \\
  Lightweight & Decrease the Armor Class of this armor by one, if possible \\
  Reinforced & Increase the Armor Class of this armor by one, if possible \\
\end{Table}



\section{Jeweling}
\textit{Jeweling} allows characters to create jewelry.

\vspace{0.4em}

If a character can craft jewelry, 
they may craft one item of jewelry whenever they level-up in a town with a Workshop.

\vspace{0.4em}

When you create jewelry, you must expend a metal ingot and any number of gemstones 
(up to the maximum number of gemstones for that jewelry type).



\section{Enchanting}
\textit{Enchanting} allows characters to imbue items with magical effects.

\vspace{0.4em}

If a character can enchant items, they may enchant one item whenever they level-up in 
a town with an Arcaneum.

\vspace{0.4em}

Enchantments count toward the maximum number of improvements allowed for a weapon or armor.

\vspace{0.4em}

To enchant a weapon or armor, you must expend a gemstone or gemstones of a type or types according to
the attributes listed in the W.I.Z.A.R.D.S. requirements for the effect that is to be added:
\begin{Table}{l c}
  \TableHeader{Attribute} & \TableHeader{Gemstone} \\
  Wisdom & Sapphire \\
  Intelligence & Opal \\
  Zest & Topaz \\
  Agility & Emerald \\
  Resilience & Garnet \\
  Dexterity & Amethyst \\
  Strength & Ruby \\
\end{Table}
\textbf{Diamonds} may be used in place of any gemstone, 
or to enchant effects that do not have a W.I.Z.A.R.D.S. requirement.

\vspace{0.4em}

Choose an effect for the enchantment:
\begin{itemize}
  \item Choose an option from the table in Appendix H.
  \item Add a perk to the item. The character equipped with the item has that perk; 
    the effect is enabled as long as they have all prerequisites, 
    even if they do not have the W.I.Z.A.R.D.S. requirements for that perk.
    The enchanted perk cannot be used to satisfy the prerequisites for any other perks.
  \item Imbue the item with a Reservoir (i.e, Life or Mana). 
    The base value for the reservoir is equal to \Roll{1d10} + 
    your character's \textbf{Intelligence} and \textbf{Zest} modifiers. 
    Reservoirs imbued in weapons \& rings may be spent to activate abilities of that weapon, 
    or of the item equipped in the hand that is equipped with the ring. 
    Reservoirs imbued in armor \& accessories may be freely spent in place of the character's reservoirs. 
    When a character's reservoirs are restored by a spell, 
    they may choose to restore the reservoir of their equipment instead.
\end{itemize}



\section{Alchemy}
\textit{Alchemy} allows characters to craft potions.

\vspace{0.4em}

If a character can do alchemy, 
they may craft enough potions to fill a number of flasks up to to their \textbf{Intelligence} modifier 
whenever they level-up in a Town with a Laboratory.

\vspace{0.4em}

Potion profiles can be found in Appendix G.


\subsection{Crafting Cost}
To craft a potion, a character requires a \textit{Bottle} (See Bottle in Chapter 2: Measurement) 
to hold the crafted potion.

\vspace{0.4em}

Then, the character must expend \textit{Ingredients} listed in the potion \textit{Recipe}.
Ingredients (e.g., herbs, minerals) determine the effects of the crafted potion,
and the ingredients listed in the recipe must be expended for each flask created.

\begin{Example}
  Say a potion's recipe lists a flask of \textit{Vinegar}, 
  a Tiny amount of \textit{Nightshade}, and a Tiny amount of \textit{Sulfur}.

  \vspace{0.4em}

  If a character is able to do alchemy, has an Intelligence modifer of +2,
  and has the opportunity to level-up in a Town with a Laboratory,
  then this character may craft 2 flasks of the potion by expending 2 empty flasks,
  2 Tiny amounts of Nightshade, and 2 Tiny amounts of Sulfur.

  \vspace{0.4em}

  If the recipe lists a flask of \textit{Vinegar} and a Tiny amount of \textit{Sulfur},
  then the same character may craft 2 flasks of the potion by expending 2 flasks of Vinegar
  and 2 Tiny amounts of Sulfur, provided the Vinegar is already stored in 2 individual flasks.
\end{Example}


\subsection{Concoctions}
If a character knows the properties of one or more materials that may serve as potion ingredients,
then they may craft a unique potion by expending 5 Tiny amounts of each of those materials,
plus 1 flask's worth of a liquid base, in order to produce one flask of potion with 5 charges.
The potion's effect is the combined effects of each of the ingredients and the liquid base.

\vspace{0.4em}

This unique combination of ingredients and liquid base becomes a recipe that may be recorded in a folio,
whereupon other characters with access to the folio may craft the same potion with the same effects
without knowing the effects of the individual materials.

\newpage

\subsection{Poisoned Weapons}
Potions can be applied to a weapon or ammunition to trigger the potion's effect on a target 
when an attack made with that weapon or ammunition deals damage to a target.

\vspace{0.4em}

To apply a potion to a weapon or ammunition, 
expend any number of charges of that potion when you level-up in a Town,
then roll Xd10, where X is the number of charges expended.

\vspace{0.4em}

If you applied the potion to ammunition, 
that result is the number of ammunition (e.g., arrows) to which the potion effect is applied.

\vspace{0.4em}

If you applied the potion to a weapon, 
that result is the number of attacks to which the potion's effect may be applied.


\subsection{Consumable Potion Effects}
\textit{Potions} are consumable items that are stored in bottles, 
and their potency and use are measured in \textit{Charges}. 

\vspace{0.4em}

Potions have an innate \textit{Consume} ability, 
whereby expending 1 charge you can activate the effect of the potion.


\subsection{Non-Consumable Potion Effects}
Some potions may have alternative effects when they are not used as a consumable.
In such a case, these potions will provide a new ability to the bottle in which they are stored.

\begin{Example}
  An explosive potion may provide the triggered ability:
  
  \vspace{0.4em}

  \textit{When this bottle is used as a projectile for a ranged attack, it shatters and is destroyed,
  dealing Xd4 damage to all characters within 1", where X is the number of charges remaining in the bottle.
  If the projectile attack hits, the radius is measured from the base of the target.
  If the projectile attack misses, draw a line from the base of the attacker through the base of the target,
  then mark a point 2" past the base of the target along that line; 
  the radius of the damage is centered on the marked point.}
\end{Example}

The crafting process for potions with non-consumable abilities is the same as for consumable potions.


